%% Starter deck for the SES-ALCF template, and the reference for every layout
%% the theme provides.
%%
%%   cp -r template my-talk && cd my-talk
%%   rename template.tex (and its three mentions in the Makefile)
%%
%% Families: a plain \begin{frame} is the navy family, \begin{seslightframe} the
%% white one; \sesalcfdark / \sesalcflight switch a whole document.
\documentclass[10pt,aspectratio=169]{beamer}
\usetheme[sectionpages]{sesalcf}

\usepackage{booktabs}
\usepackage[newfloat]{minted}
\setminted{fontsize=\scriptsize,baselinestretch=1,breaklines,bgcolor=SESpanelLight}

\title{Talk Title}
\subtitle{One line of context}
\author{Ethan Wong}
\date{\today}

\begin{document}

%% Title Slide
\begin{frame}[plain]
  \titlepage
\end{frame}

%% 1_*Section Break -- inserted by the theme's [sectionpages] option
\section{Section Break}

%% *Title, Subtitle and Bullets
\begin{frame}{Title, Subtitle and Bullets}
  \framesubtitle{The subtitle is the pptx's light-blue sentence-case line}
  \begin{itemize}
    \item First-level bullet: white text on the navy master
    \item A second bullet with \textbf{bold}, \emph{italic} and \texttt{monospace} runs
    \begin{itemize}
      \item Second-level bullet, rendered as an en dash
      \item Another second-level bullet
    \end{itemize}
  \end{itemize}
\end{frame}

%% *Title and Subtitle Only
\begin{frame}{Title and Subtitle Only}
  \framesubtitle{No body placeholder: use this for a single statement or a section lead-in}
\end{frame}

%% *Title and Subtitle: 2 column
\begin{seslightframe}
  \frametitle{Title and Subtitle: 2 Column}
  \framesubtitle{Each column gets its own heading and body}
  \begin{sescolumns}{2}
    \sescolumn{Left column}{
      \begin{itemize}
        \item Bullets, text or a table on the left
        \item Columns are equal width
      \end{itemize}}
    \sescolumn{Right column}{
      \begin{itemize}
        \item The right-hand column
        \item Same heading and body styles
      \end{itemize}}
  \end{sescolumns}
\end{seslightframe}

%% *Title and Subtitle: 3 column
\begin{seslightframe}
  \frametitle{Title and Subtitle: 3 Column}
  \framesubtitle{Swap the argument to \texttt{sescolumns} for 3, 4 or 5}
  \begin{sescolumns}{3}
    \sescolumn{First}{Text or bullets.}
    \sescolumn{Second}{Text or bullets.}
    \sescolumn{Third}{Text or bullets.}
  \end{sescolumns}
\end{seslightframe}

%% *Title and Subtitle: 4 column
\begin{seslightframe}
  \frametitle{Title and Subtitle: 4 Column}
  \begin{sescolumns}{4}
    \sescolumn{One}{Narrow columns suit short labels or numbers.}
    \sescolumn{Two}{Text or bullets.}
    \sescolumn{Three}{Text or bullets.}
    \sescolumn{Four}{Text or bullets.}
  \end{sescolumns}
\end{seslightframe}

%% *Title and Subtitle: 5 column
\begin{seslightframe}
  \frametitle{Title and Subtitle: 5 Column}
  \begin{sescolumns}{5}
    \sescolumn{One}{Short.}
    \sescolumn{Two}{Short.}
    \sescolumn{Three}{Short.}
    \sescolumn{Four}{Short.}
    \sescolumn{Five}{Short.}
  \end{sescolumns}
\end{seslightframe}

%% *Title and Subtitle: 4 block big idea
\begin{seslightframe}
  \frametitle{Title and Subtitle: 4 Block Big Idea}
  \framesubtitle{Two panels per row, each with its own heading}
  \sespanel{Big idea one}{One or two sentences of supporting text.}\hfill
  \sespanel{Big idea two}{One or two sentences of supporting text.}\par
  \vspace{0.35cm}
  \sespanel{Big idea three}{One or two sentences of supporting text.}\hfill
  \sespanel{Big idea four}{One or two sentences of supporting text.}
\end{seslightframe}

%% *Title and Subtitle: 3 column photo
\begin{seslightframe}
  \frametitle{Title and Subtitle: 3 Column Photo}
  \framesubtitle{Each column is a picture with a caption underneath}
  \begin{sesphotos}{3}
    \sesphoto{sesalcf-assets/bg-title.png}{Caption for the first picture.}
    \sesphoto{sesalcf-assets/bg-section.png}{Caption for the second.}
    \sesphoto{sesalcf-assets/bg-title.png}{Caption for the third.}
  \end{sesphotos}
\end{seslightframe}

%% Custom Layout: full-bleed media with the title on top
\begin{frame}[plain]
  \sesfullbleed[Custom Layout: Full-Bleed Media]{sesalcf-assets/bg-title.png}
\end{frame}

%% Code, in the theme's light panel
\begin{frame}[fragile]{Code}
  \framesubtitle{\texttt{minted} panels take the panel colours from the theme}
  \begin{minted}{bash}
# a shell sample
echo "hello"
  \end{minted}
\end{frame}

%% What the theme gives you
\begin{frame}{Using This Template}
  \framesubtitle{The same API works in any deck: \texttt{\textbackslash usetheme[sectionpages]\{sesalcf\}}}
  \begin{itemize}
    \item \texttt{\textbackslash title}, \texttt{\textbackslash subtitle}, \texttt{\textbackslash author}, \texttt{\textbackslash date} \(\rightarrow\) title page
    \item \texttt{\textbackslash section\{...\}} \(\rightarrow\) section page, with the \texttt{sectionpages} option
    \item \texttt{\textbackslash begin\{frame\}\{Title\}\textbackslash framesubtitle\{...\}} \(\rightarrow\) title and subtitle
    \item \texttt{seslightframe} (title via \texttt{\textbackslash frametitle}) \(\rightarrow\) the white slide family
    \item \texttt{sescolumns\{n\}}, \texttt{\textbackslash sescolumn\{head\}\{body\}} \(\rightarrow\) 2--5 column layouts
    \item \texttt{\textbackslash sespanel\{head\}\{body\}} \(\rightarrow\) big-idea panels
    \item \texttt{sesphotos\{n\}}, \texttt{\textbackslash sesphoto\{file\}\{caption\}} \(\rightarrow\) photo columns
    \item \texttt{\textbackslash sesfullbleed\{file\}} \(\rightarrow\) full-bleed media
  \end{itemize}
\end{frame}

\end{document}
